% Escreva aqui a solução da questão 4.


\textcolor{red}{ VERSÃO 1 } %ian

Seja $V$ a quantidade de anos vivida por Diofanto. Analisando o texto, as partes da vida de Diofanto indicadas são as seguintes:

``Deus lhe concedeu a \textbf{sexta parte \emph{da vida}} para ser menino'' : $\dfrac{V}{6}$.

``e uma \textbf{duodécima parte \emph{a mais}} para que os seus pêlos começassem a crescer'': $\dfrac{V}{12}$.

``Depois de uma \textbf{sétima parte \emph{adicional}}, ele se casou'': $\dfrac{V}{7}$.

``\textbf{\emph{Passados} cinco anos}, nasceu seu filho'': $+5$ anos. 

``Mas o destino decidiu que \textbf{o filho} viveria apenas \textbf{metade do tempo de \emph{vida total do pai}}'': $\dfrac{V}{2}$.

``\textbf{\emph{Após} a morte do filho}, Diofanto \textbf{viveu \emph{ainda} quatro anos}'': $+4$ anos.

%A análise consiste em perceber que: os períodos se referem sempre ao \textbf{todo} de sua vida; os períodos nunca se referem a anos "restantes"; os períodos são distintos, sem sobreposição de intervalo de tempo; os períodos narrados juntos compreendem \textbf{toda} a sua vida; a informação numérica do nascimento e morte do seu filho permite determinar numericamente quantos anos teve de vida. Fazemos essa determinação pela equação:

Disso, segue a solução da equação:

\begin{align*}
    \dfrac{V}{6} + \dfrac{V}{12} + \dfrac{V}{7} + 5 + \dfrac{V}{2} + 4 = V &\iff \dfrac{14V+7V+12V+42V}{84} + 9 = V\\
    &\iff \dfrac{14V+7V+12V+42V}{84} + 9 = V\\
    &\iff 9 = V - \dfrac{75V}{84}\\
    &\iff 9 = \dfrac{9V}{84}\\
    &\iff V = 84.
\end{align*}

Portanto, Diofanto viveu 84 anos.

\textcolor{red}{ VERSÃO 2 } %Esther

Seja a quantidade de anos vividas por Diofanto denotada por $x$. Do texto, tiramos as seguintes conclusões:

\begin{itemize}
    \item $\dfrac x6$ foi sua infância;
    \item $\dfrac x{12}$ foi sua adolescência, até os $\dfrac{x}6 + \dfrac{x}{12} = \dfrac{x}{4}$ anos;
    \item Depois de mais $\dfrac x7$ anos, se casou, aos  $\dfrac x4 + \dfrac{x}{7} = \dfrac{11x}{28}$ anos;
    \item Com mais cinco anos, teve um filho, aos  $\dfrac{11x}{28} + 5 = \dfrac{11x+140}{28}$ anos; 
    \item Após $\dfrac x2$ anos, seu filho morreu, quando Diofanto tinha $\dfrac{11x+140}{28} + \dfrac x2 = \dfrac{25x+140}{28}$ anos;
    \item Diofanto morreu aos $\dfrac{25x+140}{28} + 4 = \dfrac{25x+252}{28}$. Como aqui ele morreu, chegamos aos $x$ anos.
\end{itemize}

Disto, basta resolver a equação:

\begin{align*}
    \dfrac{25x+252}{28} = x &\iff 25x+252 = 28x\\
    &\iff 3x = 252\\
    &\iff x = 84.
\end{align*}

Então, Diofanto viveu 84 anos.







%\bigskip

%\textcolor{red}{ VERSÃO 3} %autor





\bigskip

\textcolor{red}{CHAVE DE PONTUAÇÃO:
\begin{itemize}
    \item Obter a equação escrevendo de onde vem cada expressão \emph{(1,5 ponto)};
    \item Obter a equação sem escrever de onde vem cada expressão \emph{(0,5 ponto)};
    \item Resolução da equação e resposta final  \emph{(1,5 ponto)}.
\end{itemize} }